\section{Concluzii}
\label{sec:conclusion}

Această lucrare a prezentat proiectarea și implementarea unui joc de curse auto 3D în Unity, demonstrând aplicarea practică a conceptelor de proiectare avansată a sistemelor software. Principalele contribuții ale proiectului sunt:

\begin{enumerate}
    \item \textbf{Un model de fizică vehiculară realist} bazat pe WheelCollider-e Unity, augmentat cu sisteme de control al tracțiunii, downforce simulat, asistență la aderența laterală și gestionare inteligentă a coliziunilor cu pereții. Aceste sisteme cooperează pentru a oferi o experiență de conducere stabilă și responsivă.
    
    \item \textbf{Un sistem AI funcțional} bazat pe navigarea prin waypoint-uri, cu variabilitate introdusă prin aleatorizarea poziției pe lățimea waypoint-ului, zone de boost/încetinire și mecanism de recuperare automată din blocaje.
    
    \item \textbf{Definirea traseului prin waypoint-uri} configurate în editorul Unity, soluție care unifică reprezentarea vizuală a circuitului cu datele de navigare ale AI-ului și permite forme arbitrare de traseu.

    \item \textbf{Sisteme complete de game flow} -- numărătoare inversă, contorizare tururi, pauză și ecran de final -- care formează o experiență de joc completă.
\end{enumerate}

Proiectul ilustrează importanța arhitecturii modulare în dezvoltarea de software complex: fiecare subsistem (fizică, AI, traseu, UI) este implementat ca o componentă independentă, ușor de testat și de extins.
